\section{The unrestricted all-test return criterion at every width}
\label{sec:all-width-return-corollary}

The finite obstruction of MA-PR-004 and the polynomial
criterion of MA-PR-010 combine into an exact characterization
without an initial restriction on boundary multiplicities.
All quotients, norms, translations and centering conventions
remain those of (MA-PR.1), with arbitrary fixed \(u<v\),
\(a\in\mathbb R\) and \(s\ge0\).

\paragraph{MA-PR-011 (both necessary gates for unrestricted tests).}
The following conditions are equivalent.
\begin{enumerate}
\item Every \(g\in\mathcal E\) admits a continuous
equivariant return \(S_g:Q_H\to Q_\infty\) satisfying
\[
 \begin{aligned}
          iS_g&=C_g^{Q_H},\\
          S_gU_t^{Q_H}&=e^{-at}U_t^{Q_\infty}S_g
                         \quad(t\in\mathbb R).
 \end{aligned}
 \tag{MA-PR.117}
\]
\item Every compact smooth \(g\) with \(g(0)=0\)
admits a return satisfying (MA-PR.117).
\item Every \(g\in\mathcal D_{[1,2]}\) admits a return
satisfying (MA-PR.117).
\item There is no proper nonzero retained primary factor:
\[
                  r_\rho\in\{0,m_\rho\}
                             \quad\hbox{for every }\rho.
 \tag{MA-PR.118}
\]
Once this holds, form the canonical inclusions and their
actual whole-quotient costs as in (MA-PR.101).  They satisfy
\[
       \mathcal C_n=\|\iota_n\mathscr R_n\|_{Q_H\to J_n}
                    \le C\nu_n^D
                 \quad(n\in\mathbb Z)
 \tag{MA-PR.119}
\]
for some fixed finite \(C>0,D\ge0\).
\end{enumerate}
Condition~2 is equivalently stated for all compact smooth
tests, without the value-zero restriction.
The canonical inclusion in condition~4 is invoked only
after its preceding no-proper-truncation condition is satisfied.

When these conditions hold, the quotient-valued return
is unique.  Its linear choice, jointly continuous
\(\mathcal E\)-valued interpolation lift, and common
all-order nuclear representations are exactly those
proved in MA-PR-010.  If condition~4 fails, the tests
in \(\mathcal D_{[1,2]}\) admitting no such return
contain a dense \(G_\delta\) subset.
Neither validity nor failure of either arithmetic
condition is asserted at a particular width.

\begin{proof}
The test-space inclusions give \(1\Rightarrow2\Rightarrow3\).
Suppose condition~3 holds and that there were a primary
factor with \(0<r_\rho<m_\rho\).  Put
\(\lambda=\rho-a\), choose
\(\eta\in C_c^\infty((1,2))\) with integral one, and set
\[
              g(x)=e^{-\lambda x}\eta(x),
              \qquad g(0)=0,\qquad
              \mathcal Mg(\lambda)=1.
 \tag{MA-PR.120}
\]
This is an actual test in \(\mathcal D_{[1,2]}\).
The finite module necessity in MA-PR-004,
(MA-PR.25)--(MA-PR.29), requires
\(\mathcal Mg(\lambda)=0\) for every return at that
proper factor, since \(\min(m_\rho-r_\rho,r_\rho)\ge1\).
It applies to any global return by the actual finite
Hilbert block and full smooth jet restrictions proved
there.  Equation (MA-PR.120) contradicts that necessity.
Thus condition~3 forces (MA-PR.118).

There are now no eligibility restrictions.  Apply the
\(4\Rightarrow1\) implication of MA-PR-010 to the
entire fixed-support test space in the present condition~3.
In particular its proved Baire bound (MA-PR.114) and
uniform inverse-class argument (MA-PR.115) give
(MA-PR.119).  This proves \(3\Rightarrow4\).
Conversely, (MA-PR.118)--(MA-PR.119) are precisely the
hypotheses of the \(1\Rightarrow2\) implication of
MA-PR-010.  Its actual series (MA-PR.108)--(MA-PR.109)
constructs the return for every \(g\in\mathcal E\),
proving \(4\Rightarrow1\) and the stated uniqueness,
topology and nuclear properties.
The version of condition~2 without \(g(0)=0\) follows
from the same inclusions between conditions~1 and~3.

For the conditional failure statement, first suppose
there is a proper retained factor as above.  The functional
\(g\mapsto\mathcal Mg(\lambda)\) is continuous on
\(\mathcal D_{[1,2]}\), and (MA-PR.120) shows it is
nonzero.  Therefore
\[
       \{g\in\mathcal D_{[1,2]}:
                         \mathcal Mg(\lambda)\ne0\}
 \tag{MA-PR.121}
\]
is open and dense.  Density follows also by adding an
arbitrarily small nonzero multiple of the test in
(MA-PR.120) to any element of its kernel.
Every test in (MA-PR.121) fails the finite return
condition.  This open dense set is itself a
\(G_\delta\).  If instead (MA-PR.118) holds but the
polynomial estimate fails, the dense \(G_\delta\)
alternative (MA-PR.104) applies.  These cases exhaust
failure of condition~4 and prove the assertion.
\end{proof}

The first gate is algebraic at the actual boundary
jets; the second is a measured analytic growth condition
on the whole Hilbert quotient.  For \(s\le1\) only
the first gate is automatic.  At higher weights,
absence of boundary zeros also removes only that
first obstruction.  No polynomial estimate is inferred
from these facts alone.

This characterizes the actual quotient returns in
(MA-PR.117).  It does not erase the discarded nuclear
correction of (MA-PR.116), identify an original Meyer
common-space projection on a new Hilbert domain, or
prove an arithmetic zero-location or positivity theorem.
The exact proof inputs are the complete local module
obstruction MA-PR-004 and the complete norm/topology
equivalence MA-PR-010; no new outside source is used.
